% !TEX root = ../eigenfaces_algorithm.tex
%=====================================================
\section{关键困难与 Gram 技巧：在 280 那一侧算}

\subsection{困难：矩阵太大}

$\mat{C}=\frac{1}{n-1}\mat{X}\mT\mat{X}\in\R^{4096\times4096}$。
对它做完整特征分解约需 $O(d^3)\approx 7\times10^{10}$ 次运算，
光把它存下来就要 $4096^2\times 8\text{B}\approx 134$ MB。
而每个特征向量才用一次——太奢侈了。

\subsection{Gram 技巧：三个等式解决问题}

改看另一个矩阵：
\[
\mat{G}=\mat{X}\mat{X}\mT\in\R^{280\times 280}.
\]
它小得多。它的元素有直接含义：$G_{ij}=\inner{\mathbf{x}_i}{\mathbf{x}_j}$
是第 $i$ 张残差脸和第 $j$ 张残差脸的\textbf{相似度}——所以 $\mat{G}$ 叫 Gram 矩阵（相似度矩阵）。

\textbf{断言：}若 $\mat{G}\mathbf{u}=\lambda\mathbf{u}$ 且 $\lambda>0$，
则 $\mathbf{v}=\mat{X}\mT\mathbf{u}$ 是 $\mat{X}\mT\mat{X}$ 的特征向量，特征值同为 $\lambda$。

\textbf{证明（两行）：}
\[
\mat{X}\mT\mat{X}\,(\mat{X}\mT\mathbf{u})
=\mat{X}\mT\,(\mat{X}\mat{X}\mT\,\mathbf{u})
=\mat{X}\mT(\lambda\mathbf{u})
=\lambda\,(\mat{X}\mT\mathbf{u}).\qquad\blacksquare
\]

\textbf{归一化：}$\mat{X}\mT\mathbf{u}$ 不是单位向量，需要除以它的长度：
\[
\norm{\mat{X}\mT\mathbf{u}}^2
=\mathbf{u}\mT\mat{X}\mat{X}\mT\mathbf{u}
=\lambda\,\mathbf{u}\mT\mathbf{u}=\lambda
\quad\Longrightarrow\quad
\boxed{\ \mathbf{v}=\frac{\mat{X}\mT\mathbf{u}}{\sqrt{\lambda}}\ },\qquad \norm{\mathbf{v}}=1 .
\]

\subsection{两套谱的关系与“空房间”}

$\mat{X}\mT\mat{X}$ 与 $\mat{X}\mat{X}\mT$ 的非零特征值\textbf{完全相同}：
$\mat{X}$ 的秩 $\le 279$，所以两边都至多 279 个非零特征值。
协方差版本只差常数：$\mat{C}$ 的特征值为 $\lambda_i/(n-1)$。
此外若 $\mathbf{u}_i\perp\mathbf{u}_j$（Gram 侧），则
$\inner{\mathbf{v}_i}{\mathbf{v}_j}=\tfrac{1}{\sqrt{\lambda_i\lambda_j}}\mathbf{u}_i\mT\mat{X}\mat{X}\mT\mathbf{u}_j
=\tfrac{\lambda_j}{\sqrt{\lambda_i\lambda_j}}\inner{\mathbf{u}_i}{\mathbf{u}_j}=0$：
\textbf{特征脸彼此也是正交的}——正交性在小矩阵那侧算完就自动继承到大矩阵这侧。

\begin{intuition}
一个直观比喻：$\mat{X}$ 是一座桥，左岸住着 280 个“人”（$\mathbf{u}$ 们），
右岸住着 4096 个“像素方向”（$\mathbf{v}$ 们）。
桥上真正通行的“变异模式”只有 $\le 279$ 条。
在左岸（人这一侧）把这条条模式找清楚，桥的对称性会自动告诉我们右岸的对应者是谁：
$\mathbf{v}=\mat{X}\mT\mathbf{u}/\sqrt{\lambda}$。
这就是“$n\ll d$ 时先算 $n\times n$ 那一侧”的全部秘密。
\end{intuition}

\subsection{与 SVD 的对应（一张对照表）}

瘦奇异值分解 $\mat{X}=\mat{U}\mat{\Sigma}\mat{V}\mT$ 中
$\sigma_i=\sqrt{\lambda_i}$，于是：
\[
\text{Gram 侧：}\ \mat{G}=\mat{U}\mat{\Lambda}\mat{U}\mT,\qquad
\text{特征脸：}\ \mat{V}=\mat{X}\mT\mat{U}\mat{\Lambda}^{-1/2}.
\]
一句话总结本题的算法本质：\textbf{“特征脸 = 只用 $n$ 侧信息求出 $\mat{X}$ 的右奇异向量。”}
